% ============================================================================
% RELATED WORK (REVISED - Collaborative Learning Frame)
% ============================================================================

\section{Related Work}
\label{sec:related}

Our work builds upon and learns from three research areas: adaptive model selection, cost-aware inference optimization, and transfer learning for decision-making systems. We position \ours{} as a complementary alternative that prioritizes accessibility and ease of deployment.

\subsection{LLM Routing and Model Selection}

\paragraph{Cascading Systems.} FrugalGPT~\cite{chen2023frugalgpt} pioneered cost-optimized LLM routing through sequential model invocation with learned verification thresholds, demonstrating impressive cost-quality trade-offs on fixed model pools. RouteLLM~\cite{ong2024routellm} extends this paradigm by training preference-based classifiers to route between two models. Both systems validate that \emph{adaptive model selection works}---a foundational insight upon which we build.

However, these approaches face operational barriers that limit accessibility:

\begin{table}[t]
\centering
\footnotesize
\caption{\textbf{Operational Requirements Comparison.} We learn from prior systems' strengths while reducing deployment barriers.}
\label{tab:operational_comparison}
\begin{tabular}{lccc}
\toprule
\textbf{Requirement} & \textbf{FrugalGPT} & \textbf{RouteLLM} & \textbf{Ours} \\
\midrule
Calibration Data & hundreds to thousands & thousands of pairs & <100 examples \\
Domain Expertise & High (scorer design) & Medium (classifier) & Low (metadata) \\
Model Updates & Full recalibration & Retrain classifier & Register + explore \\
Setup Time & Days & Hours & Minutes \\
\bottomrule
\end{tabular}
\end{table}

\paragraph{The Calibration Barrier.} FrugalGPT requires users to provide labeled calibration datasets (often hundreds to thousands of examples) and train domain-specific scoring functions to verify output quality~\cite{chen2023frugalgpt}. This setup cost is prohibitive for students, independent researchers, and small teams lacking ML expertise or annotated data. Our approach eliminates this barrier through \emph{shippable priors}---pre-trained covariance matrices that provide immediate routing intelligence without user-provided calibration data.

\paragraph{The Expertise Barrier.} Effective cascade design requires deep knowledge of LLM capabilities to manually curate model chains and set verification thresholds. A user must understand which models form appropriate "tiers" and how to balance precision-recall trade-offs in scoring functions. \ours{} addresses this through autonomous online learning: the system discovers model strengths through production traffic, removing the need for upfront architectural decisions.

\paragraph{The Maintenance Barrier.} Static systems require expensive recalibration when model capabilities evolve (e.g., monthly GPT-4o updates) or when new models launch. In rapidly evolving markets where cheaper alternatives emerge weekly, $O(N)$ maintenance becomes a bottleneck. Our decoupled architecture (\S\ref{sec:scalability}) enables $O(1)$ model onboarding: new models are registered via public metadata and evaluated automatically through online exploration.

\paragraph{Learning from Cascades.} FrugalGPT's cascading architecture excels at achieving high reliability through sequential verification---a strength we incorporate in our \emph{Hybrid Mode} (\S\ref{sec:hybrid}). However, we address its accessibility limitations by using prediction as a gating mechanism: the bandit performs ex-ante risk assessment to decide \emph{when} verification is needed, rather than requiring users to manually configure fixed chains. This preserves FrugalGPT's reliability while enabling single-pass routing over 80+ models without latency penalties.

\paragraph{A Critical Limitation: Confident Failure.} Through systematic evaluation (\S\ref{sec:confident_failure}), we identify a failure mode in reactive verification systems: scoring functions incorrectly validate 35\% of erroneous outputs from cost-optimized models when outputs are plausible but semantically incorrect. Our predictive architecture mitigates this by routing high-uncertainty prompts directly to frontier models before generation, preemptively avoiding the pathway where confident failures occur. This insight suggests that ex-ante risk assessment complements ex-post verification rather than replacing it.

\subsection{Multi-Armed Bandits and Online Learning}

\paragraph{Contextual Bandits.} Contextual multi-armed bandits~\cite{li2010contextual,chu2011contextual} provide a principled framework for sequential decision-making under uncertainty, with LinUCB~\cite{li2010contextual} achieving provably optimal regret bounds under linear reward assumptions. However, standard bandits suffer from cold-start inefficiency: convergence requires hundreds to thousands of samples, translating to substantial exploration costs (\$50--200) that price out resource-constrained users.

\paragraph{Our Contribution: Warm-Start Mechanism.} We address cold-start inefficiency through offline expert distillation (\S\ref{sec:distillation}). By pre-training covariance matrices on teacher supervision, we reduce day-one regret by 64.6\% (\S\ref{sec:rq1}) while preserving the bandit's adaptive plasticity. This differs from standard transfer learning~\cite{taylor2009transfer} by explicitly encoding uncertainty structure (covariance) rather than point estimates alone, enabling informed exploration from the first query.

\paragraph{Plasticity Under Drift.} A key advantage of online learning over static systems is autonomous adaptation to concept drift. We validate this through poisoned prior experiments (\S\ref{sec:rq2}), demonstrating that \ours{} successfully corrects teacher bias and model capability changes without manual intervention. This plasticity ensures long-term reliability as the model ecosystem evolves.

\subsection{LLM Benchmarking and Evaluation}

\paragraph{Static Benchmarks.} Comprehensive evaluation frameworks like HELM~\cite{liang2022holistic}, BIG-bench~\cite{srivastava2022beyond}, and LMSYS Chatbot Arena~\cite{zheng2023judging} provide valuable model comparisons. We leverage these benchmarks for offline prior distillation and external validation (\S\ref{sec:rq1}).

\paragraph{The Benchmark Trap.} However, initializing routing policies from public benchmark leaderboards yields worse performance than random initialization (Table~\ref{tab:benchmark_trap}). This occurs because benchmark-optimal models diverge systematically from deployment-optimal models: benchmarks evaluate broad capability across standardized distributions, whereas effective routing requires utility optimization under deployment-specific constraints (cost, latency, local task distributions). This finding suggests that routing systems require domain-specific calibration rather than global rankings.

\subsection{Design Philosophy: Complementary Alternatives}

Rather than positioning \ours{} as a replacement for existing systems, we view it as a complementary alternative optimized for different operational contexts:

\begin{itemize}[leftmargin=*,nosep]
    \item \textbf{FrugalGPT excels when:} Users have labeled calibration data, domain expertise to design scorers, and stable prompt distributions where upfront optimization costs amortize over long deployments.
    \item \textbf{BanditGPT excels when:} Users lack calibration data or ML expertise, face rapidly evolving model ecosystems, or prioritize immediate deployment over asymptotic optimality.
\end{itemize}

Our core insight is that \emph{accessibility requires ease of use, not just cost reduction}. By eliminating calibration requirements, reducing setup time from days to minutes, and enabling autonomous adaptation, we expand the frontier of who can deploy adaptive routing---from ML experts to students, researchers, and startups without specialized infrastructure.

\subsection{Positioning Summary}

\begin{table}[t]
\centering
\footnotesize
\caption{\textbf{Learning from Prior Work.} \ours{} combines strengths while addressing accessibility barriers.}
\label{tab:related_comparison}
\begin{tabular}{p{2.2cm}p{2.4cm}p{2.4cm}}
\toprule
\textbf{System} & \textbf{Strength (We Learn)} & \textbf{Limitation (We Address)} \\
\midrule
FrugalGPT & High reliability via cascading & Heavy setup; manual chains \\
RouteLLM & Preference learning & Static; requires retraining \\
Std.\ Bandits & Principled exploration & Cold-start (\$50--200) \\
Leaderboards & Global capability rankings & Cost-agnostic; not contextual \\
\bottomrule
\end{tabular}
\end{table}

The fundamental innovation enabling \ours{}'s accessibility is the decoupling of \emph{prompt difficulty learning} (performed once offline, shipped as priors) from \emph{model capability profiling} (performed continuously online). This architectural separation allows users to deploy immediately without calibration data, while the system autonomously refines routing decisions through production traffic. By reducing both economic barriers (61\% cost reduction) and expertise barriers (minutes vs.\ days of setup), we aim to expand adaptive routing from specialized ML teams to the broader community of students, researchers, and practitioners.

